% Werke X/1 (1917), native206–209, printed199–202.
\begingroup
\setlength{\headheight}{25pt}
\makeatletter\tagsleft@true\let\veqno\@@leqno\makeatother
\SourceLeaf{206}\ArithmeticHead{199}{R}{}
\fancyhead[C]{\scriptsize THEORIA SINUS LEMNISCATICI UNIVERSALISSIME ACCEPTI.}
\NumberedItem{7}
\noindent [S.\ 37]
\[\operatorname{sl}n\varpi=\operatorname{t[ang]}u\pi\]
\[\frac\pi\varpi du=\operatorname{cl}n\varpi dn\]
\[\operatorname{cl}n\varpi=\frac\pi\varpi\left(\frac4{e^{\frac12\pi}+e^{-\frac12\pi}}\operatorname{c[o]s}n\pi+\ldots\right)\]
\[\pi u=\frac4{e^{\frac12\pi}+e^{-\frac12\pi}}\sin n\pi
+\frac13.\frac4{e^{\frac32\pi}+e^{-\frac32\pi}}\sin3n\pi+\text{etc.}\]
\par\hfill Dem[onstratum.]
\[1+2p\cos\varphi+2pp\cos2\varphi+\text{etc.}=\frac{\frac1p-p}{\frac1p+p-2\cos\varphi}\]
\[1-2p\cos\varphi+2pp\cos2\varphi-\text{etc.}=\frac{\frac1p-p}{\frac1p+p+2\cos\varphi}\]
\[p\cos\varphi+p^3\cos3\varphi+\text{etc.}=\frac{\left(\frac1p-p\right)\cos\varphi}{\left(\frac1p+p\right)^2-4\cos\varphi^2}\]
\[2p\sin\varphi+\frac23p^3\sin3\varphi+\text{etc.}=\operatorname{Arctg}\frac{2\sin\varphi}{\frac1p-p}\]
\[\frac12\pi u=\operatorname{Arctg}\frac{2\sin\varphi}{e^{\frac12\pi}-e^{-\frac12\pi}}
-\operatorname{Arctg}\frac{2\sin\varphi}{e^{\frac32\pi}-e^{-\frac32\pi}}+\text{etc.}\]
\ClosingRule
\clearpage
\SourceLeaf{207}\ArithmeticHead{200}{L}{ANALYSIS.\quad NACHLASS.}
\NumberedItem{8}
\noindent [S.\ 37]
\begin{center}\textbf{Problema.}\end{center}
\par\indent Summare seriem
\[\begin{aligned}
&\left(\frac1{x+\frac1x}\right)^2+\left(\frac1{x^3+\frac1{x^3}}\right)^2
+\left(\frac1{x^5+\frac1{x^5}}\right)^2\ldots\\
&\qquad=\frac{\dfrac1{xx}+\dfrac4{x^8}+\dfrac9{x^{18}}+\dfrac{16}{x^{32}}\ldots}
{1+\dfrac2{xx}+\dfrac2{x^8}+\dfrac2{x^{18}}\ldots}
\end{aligned}\]
\ClosingRule
\noindent [S.\ 38]
\[{}'S\varphi=\frac1{GS\varphi+H\frac1{S\varphi}}\qquad\sin v=\operatorname{tg}\frac12{}'v^2\]
\begin{equation*}\tag*{1)}G+H=1\end{equation*}
\begin{equation*}\tag*{2)}
\begin{aligned}
&\surd\left(\frac1{\sin{}'v\operatorname{tang}{}'v}-\frac1{\operatorname{tang}{}'v}\right)
\times\surd\left(\frac1{\sin v\operatorname{tg}v}-\frac1{\operatorname{tg}v}\right)\times\\
&\qquad\left(G+H\frac1{\dfrac1{\sin v\operatorname{tg}v}-\dfrac1{\operatorname{tg}v}}\right)=1
\end{aligned}
\end{equation*}
\[\frac{1+\cos{}'v}2=H\]
\[\frac{1-\cos{}'v}2=G.\]
\ClosingRule
\noindent [S.\ 39]
\[\begin{aligned}
\operatorname{sl}\left(\tfrac12\varpi+(\tfrac12+k)\varpi'\varepsilon\right)
&=-\operatorname{cl}(\tfrac12+k)\varpi'\varepsilon\\
&=-\frac1{c'(\frac12+k)\varpi'}=\frac1{s'k\varpi'}
\end{aligned}\]
\[\frac1{\operatorname{sl}\left(\frac12\varpi+(\frac12+k)\varpi'\varepsilon\right)+u}=s'(k\varpi'-u\varepsilon)\]
\[\frac{\partial s}{\partial\varphi}=\surd(1-xx)\surd\left(1+\frac{xx}{\mu\mu}\right)\]
\clearpage
\SourceLeaf{208}\ArithmeticHead{201}{R}{}
\fancyhead[C]{\scriptsize THEORIA SINUS LEMNISCATICI UNIVERSALISSIME ACCEPTI.}
{\small
\[\begin{array}{c|c|c}
\frac12\varpi+(\frac12+k)\varpi'\varepsilon&0&
\frac12\left(e^{(\frac12+k)\frac{\varpi'}\varpi\pi}+e^{-(\frac12+k)\frac{\varpi'}\varpi\pi}\right)\\[1.2em]
\frac12\varpi+(\frac12+k)\varpi'\varepsilon+u&\mp u\varepsilon&
.\quad.\quad.\quad.\quad.\quad.\qquad\qquad+\nu
\end{array}\]}
\ClosingRule
\[\begin{aligned}
(P\varphi)^2={}&\frac{\surd\frac\pi\varpi}{2^{\frac74}\cos\frac18\pi}
\left(1+2e^{-2\pi}\cos4\varphi+2e^{-8\pi}\cos8\varphi+\text{etc.}\right)\\
&-\frac{\surd\frac\pi\varpi}{2^{\frac74}\sin\frac18\pi}
\left(2e^{-\frac12\pi}\cos2\varphi+2e^{-\frac92\pi}\cos6\varphi
+2e^{-\frac{25}2\pi}\cos10\varphi+\text{etc.}\right)
\end{aligned}\]
\[\begin{aligned}
(Q\varphi)^2={}&2^{\frac34}\cos\frac18\pi\surd\frac\pi\varpi
\left(1+2e^{-2\pi}\cos4\varphi+\ldots\right)\\
&+2^{\frac34}\sin\frac18\pi\surd\frac\pi\varpi
\left(2e^{-\frac12\pi}\cos2\varphi+\text{etc.}\right)
\end{aligned}\]
\NumberedItem{9}
\noindent [S.\ 40]
\par\indent Series
\[1-x-x^3+x^6+x^{10}-\text{etc.}\]
\noindent produci videtur ex evolutione producti infiniti
{\small
\[(1-x)(1-x^3)(1-x^5)(1-x^7)(1-x^9)\ldots(1-x^4)(1-x^8)(1-x^{12})\ldots\]}
\par\indent Ponendo
\[(1-z)(1-zz)(1-z^3)\ldots=\Pi z\]
\noindent erit
\begin{equation*}\tag*{1)}\frac{\Pi z\Pi z^4}{\Pi z^2}=1-z-z^3+z^6+z^{10}\ldots\end{equation*}
\begin{equation*}\tag*{2)}\frac{\Pi z\Pi z}{\Pi z^2}=1-2z+2z^4-2z^9+2z^{16}-\text{etc.}\end{equation*}
\begin{equation*}\tag*{3)}\frac{(\Pi zz)^2}{\Pi z}=1+z+z^3+[z^6+z^{10}+\ldots]\end{equation*}
{\small
\[\begin{aligned}
(\Pi z)^3={}&(1-z-z^3+z^6+z^{10}\ldots)(1-2z+2z^4-2z^9+2z^{16}\ldots)\\
&\qquad(1-2z^2+2z^8-2z^{18}\ldots)\\
={}&(1+z+z^3+z^6\ldots)(1-2z+2z^4\ldots)^2
\end{aligned}\]}
\par\noindent X1.\hfill26
\clearpage
\SourceLeaf{209}\ArithmeticHead{202}{L}{ANALYSIS.\quad NACHLASS.}
\[p=\frac{(\Pi(-z^4))^2}{\Pi z^8};\quad
q=\frac{2z\Pi(-z^8)\Pi z^{32}}{\Pi z^{16}}=\frac{2z(\Pi z^{16})^2}{\Pi z^8}\]
\[\begin{array}{c|c}
\begin{aligned}
p-q&=\frac{\Pi z.\Pi z}{\Pi z^2}\\[.8em]
p+q&=\frac{\Pi(-z)\Pi(-z)}{\Pi zz}\\[.8em]
\frac{\Pi z\Pi(-z)}{\Pi zz}&=\frac{\Pi zz\Pi zz}{\Pi z^4}
\end{aligned}
&\begin{aligned}
q&=\frac{2z\Pi z^{16}\Pi z^{16}}{\Pi z^8}\\[.8em]
p&=\frac{(\Pi z^8)^5}{(\Pi z^4)^2(\Pi z^{16})^2}\\[.8em]
p-q&=\frac{\Pi z\Pi z}{\Pi zz}\\[.8em]
p+q&=\frac{(\Pi zz)^5}{(\Pi z)^2(\Pi z^4)^2}
\end{aligned}
\end{array}\]
\[\frac1{M\cos\varphi}=\frac{(\Pi z^4)^6}{(\Pi zz)^4}\quad\text{\textit{Demonstr.}}\]
\NumberedItem{10}
\noindent [S.\ 41]
\par\indent Si
\[\begin{aligned}
&(1-a^2z^2)(1-a^2z^4)(1-a^2z^6)\ldots\left(1-\frac1{a^2}z^2\right)\left(1-\frac1{a^2}z^4\right)\ldots\\
&\hspace{15em}\times\left(a-\frac1a\right)
\end{aligned}\]
\noindent producere supponitur
\[P\left(a-\frac1a\right)+Q\left(a^3-\frac1{a^3}\right)+R\left(a^5-\frac1{a^5}\right)\ldots,\]
\noindent producet
{\small
\[\begin{gathered}
(1-a^2z^4)(1-a^2z^6)\ldots\times\left(1-\frac1{aa}\right)\left(1-\frac1{a^2}zz\right)\ldots\times\left(az-\frac1{az}\right)\\
P\left(az-\frac1{az}\right)+Q\left(a^3z^3-\frac1{a^3z^3}\right)+\text{etc.}
\end{gathered}\]}
\noindent Quare
\[\begin{aligned}
&\frac{1-aazz}{1-\frac1{aa}}\left(P\left(az-\frac1{az}\right)+Q\left(a^3z^3-\frac1{a^3z^3}\right)+\text{etc.}\right)
\times\frac{a-\frac1a}{az-\frac1{az}}\\
&\qquad=P\left(a-\frac1a\right)+Q\left(a^3-\frac1{a^3}\right)\ldots
\end{aligned}\]
\noindent Hinc sequitur
\[Q=zzP,\quad R=z^4Q,\quad S=z^6R\text{ etc.}\]
\ClosingRule
\clearpage\endgroup
